\section{The primitive incidence polarization in every characteristic}
\label{sec:T4}

All varieties in this section are over an algebraically closed field $k$.
We work with rational correspondences and with $\ell$-adic cohomology,
$\ell\ne\operatorname{char}k$; Tate twists in displayed K\"unneth components
are suppressed when no confusion is possible. A rational polarization means
an element of $\operatorname{NS}(B)_{\mathbf Q}$ a positive integral multiple
of which is the class of an ample line bundle. It does not mean a principal
polarization or an integral universal cycle.

Write $P^p(X)=\operatorname{Pic}^p(X)$ for Saito's higher Picard variety:
\[
 A^p(X)/K^p(X)\simeq P^p(X)(k),\qquad
 A^p(X)=\operatorname{CH}^p(X)_{\rm alg}.
\]
Here $K^p$ is the incidence kernel, tested by divisor-valued
correspondences. The construction, rational Poincar\'e cycles and the
canonical duality isogeny are the cited results
\cite[Theorem~3.6, \S4.1, Proposition~4.4]{Saito1979}.
The notation is the precise incidence interpretation of cote~44,
scan~2, Theorems~1--3, and scan~8. It must not be replaced without proof
by a universal regular algebraic representative.

\begin{lemma}[JAC-T4.1: primitive incidence summand]
\label{lem:primitive-summand}
Let $A$ be an abelian variety of dimension $n=2m-1$, and let $h$ be a
rational ample divisor class. Let $e_h$ denote the cohomological projector
of $H^n(A)$ onto
\[
 H^n(A)_{\rm prim}=\ker\bigl(L_h:H^n(A)\longrightarrow H^{n+2}(A)\bigr)
\]
along $L_hH^{n-2}(A)$, extended by zero in other degrees.
It is represented by a rational algebraic correspondence. It induces an
idempotent on $P^m(A)$ in the isogeny category. Its image $B=P^m(A)_{\rm prim}$
has a nondegenerate symmetric rational incidence correspondence
$d_B\in\operatorname{NS}(B\times B)_{\mathbf Q}$, rigidified along the two
axes. For a rational Poincar\'e family $z\in\operatorname{CH}^m(B\times A)_{\mathbf Q}$
inducing the inclusion $B\longrightarrow P^m(A)$, it is
\begin{equation}\label{eq:incidence}
 d_B=\operatorname{pr}_{12*}(z_{13}\cdot z_{23})_{11}.
\end{equation}
The subscript $11$ removes the pullbacks of divisor classes on the two
factors. The family may be chosen with cohomology entirely in
$H^1(B)\otimes H^n(A)_{\rm prim}$.
\end{lemma}

\begin{proof}
The Lefschetz standard conjecture and the algebraic K\"unneth projectors
hold for abelian varieties; for an exact classical locator see
\cite[Theorem~2A11]{Kleiman68}. Thus the primitive projector is algebraic.
The incidence-detection lemma Lemma~\ref{jac:P01} says that a correspondence
whose relevant odd cohomological action is zero acts as zero on higher
Picard varieties. In particular, its action depends only on the
homological class of that component. Applying this to $e_h^2-e_h$ gives
an actual idempotent in $\operatorname{End}^0(P^m(A))$.

The Poincar\'e pairing on $H^n(A)$ makes its primitive summand orthogonal
to $L_hH^{n-2}(A)$: if $L_hu=0$, then
$\int_Au\,L_hv=\int_A(L_hu)\,v=0$.
Hard Lefschetz makes this an orthogonal direct sum of nondegenerate
spaces. Consequently $e_h$ is self-transpose in degree $n$. Saito's
transpose compatibility for incidence duality
\cite[Proposition~4.4, especially (4.4.1)--(4.4.2)]{Saito1979}, together
with incidence detection, says that the induced idempotent is
self-adjoint for the canonical rational duality of $P^m(A)$.
Writing the resulting isogeny decomposition as
$P^m(A)\sim B\times B'$, this duality is block diagonal. Since the full
duality is an isogeny, each block is an isogeny. In particular the block
on $B$ is nondegenerate.

Choose Saito's Poincar\'e cycle and divide it by its positive integer
$k_m$ so that it induces the identity, then pull it back to $B\times A$.
Apply the degree-one K\"unneth projector on the parameter factor and
$e_h$ on the $A$ factor. These operations preserve the required map to
the incidence quotient and give the asserted mixed cohomological
component. Subtracting constant restrictions rigidifies the family at
the origin without changing this component. Composition of
correspondences gives \eqref{eq:incidence}. Equivalently, the
unnormalized universal cycle would multiply this correspondence by
$k_m^2$, a positive rational square. Symmetry follows by interchanging
the two parameter factors. The block argument proves its
nondegeneracy. This also proves independence of the chosen normalized
family, in the N\'eron--Severi group, from Saito's incidence identities.
\end{proof}

\begin{lemma}[JAC-T4.2: transfer of the curve test to an abelian product]
\label{lem:curve-test}
Let $C$ be a smooth projective pointed curve of positive genus $b$, and
let $A$ have dimension $2m-1$. Suppose a rational algebraic class
$\xi\in H^1(C)\otimes H^{2m-1}(A)$ satisfies
$(1\otimes L_h)\xi=0$. Let $J=J(C)$, let
$a:C\longrightarrow J$ be its Abel map, and let $\theta$ be its canonical
principal polarization. There is a rational algebraic class
$w\in H^1(J)\otimes H^{2m-1}(A)$ with
$(a\times1)^*w=\xi$. For $H=\theta+h$ on $J\times A$,
\begin{equation}\label{eq:curve-transfer}
 H^bw=0,\qquad
 \int_{J\times A}w^2H^{b-1}
 =(b-1)!\int_{C\times A}\xi^2.
\end{equation}
Hence 032, Corollary~8.4, gives
\[
 (-1)^m\int_{C\times A}\xi^2\ge0,
\]
with equality if and only if $\xi=0$.
\end{lemma}

\begin{proof}
The Abel map identifies $h^1(C)$ and $h^1(J)$ by rational algebraic
correspondences. One can use its graph, the divisor correspondence of
the Poincar\'e bundle, and the inverse Lefschetz correspondence on $J$
to obtain an algebraic inverse to $a^*:H^1(J)\to H^1(C)$.
More explicitly, $w=(b-1)!\,(L_\theta^{b-1})^{-1}(a\times1)_*\xi$ on the indicated mixed component; the inverse is algebraic by the Lefschetz standard conjecture for $J$. Indeed, $a_*a^*x=x\theta^{b-1}/(b-1)!$ for $x\in H^1(J)$. Apply this inverse to the first factor of $\xi$ and project to
$H^1(J)\otimes H^{2m-1}(A)$; the latter projector is algebraic because
both varieties in the product are abelian. This constructs $w$ as an
algebraic class, not merely as a cohomological tensor. It is unique as a
cohomological tensor, and $h\,w=0$.

The dimension of $J\times A$ is $b+2m-1$, and $w$ has codimension $m$.
Every term in $(\theta+h)^bw$ with a positive power of $h$ vanishes.
The remaining term $\theta^bw$ also vanishes because its degree on $J$
is $2b+1$. Thus $w$ is primitive for the product ample class.
Likewise $w^2H^{b-1}=w^2\theta^{b-1}$.
The canonical polarization of the Jacobian gives
\[
 a_*[C]=\frac{\theta^{b-1}}{(b-1)!}
 \quad\hbox{in cohomology}.
\]
Only its pairing with $H^2(J)$ is needed here. In this form it follows
immediately from the identification of the canonical alternating
pairings on $H^1(C)$ and $H^1(J)$: in a symplectic basis
$\theta=\sum_{i=1}^b x_i\wedge y_i$, expansion of
$\theta^{b-1}$ gives $(b-1)!$ times the complementary wedges.
The identification of the pairings is Jacobian autoduality
\cite[Chapter~III, Theorem~6.8 and Summary~6.11]{MilneAV}.
Projection formula now gives \eqref{eq:curve-transfer}.

Apply the Hodge standard conjecture in the exact form of 032,
Corollary~8.4, to the abelian variety $J\times A$, codimension $m$ and
ample class $H$. Its degree exponent is
$(b+2m-1)-2m=b-1$, and its sign is $(-1)^m$.
The factor $(b-1)!$ is positive. The equality assertion follows from
positive definiteness and the injectivity of $a^*$ in degree one.
\end{proof}

\begin{theorem}[JAC-T4.3: cote 44, Theorem 4, on the primitive incidence quotient]
\label{thm:T4}
Accept 032, Corollary~8.4, as in the task specification. For every
abelian variety $A$ of odd dimension $2m-1$ over an algebraically closed
field of any characteristic, and every rational ample class $h$, the
symmetric homomorphism on $P^m(A)_{\rm prim}$ defined by
\[
 (-1)^m d_B
\]
is a rational polarization. Equivalently, with
$\phi_L(x)=t_x^*L\otimes L^{-1}$, the rational divisor class
\begin{equation}\label{eq:T4-divisor}
 \ell_B=\frac{(-1)^m}{2}\,\Delta_B^*d_B
 \quad\in\operatorname{NS}(B)_{\mathbf Q}
\end{equation}
is ample. After multiplication by a positive integer it defines an
actual polarization in every characteristic, including characteristic
$2$.
\end{theorem}

\begin{proof}
Use Lemma~\ref{lem:primitive-summand} to choose a normalized universal
family $z$ and the nondegenerate symmetric incidence correspondence
$d_B$. For a symmetric rigidified divisorial correspondence, the
class $\Delta^*d_B/2$ is the rational N\'eron--Severi class inducing
that correspondence. Indeed, for a line bundle $L$ the associated
biextension class is
$m_B^*c_1(L)-p_1^*c_1(L)-p_2^*c_1(L)$, whose diagonal pullback is
$[2]^*c_1(L)-2c_1(L)=2c_1(L)$.
Thus \eqref{eq:T4-divisor} induces $(-1)^md_B$, which is still
nondegenerate.

Let $C\to B$ be the normalization of any integral curve in $B$.
It is nonconstant, so $C$ has positive genus. Translate it so that a
chosen point maps to the origin; translations do not change
N\'eron--Severi classes. Pulling back $z$ gives $z_C$ on $C\times A$,
whose cohomology is in
$H^1(C)\otimes H^{2m-1}(A)_{\rm prim}$.
Refined pullback and projection formula for the incidence product give
\begin{equation}\label{eq:degree-test}
 \deg_C(\ell_B)=\frac{(-1)^m}{2}
       \deg_{C\times A}(z_C\cdot z_C).
\end{equation}
There are no omitted degree-zero or degree-two parameter components:
the family was projected to the degree-one component before this
calculation. Lemma~\ref{lem:curve-test} makes the right side
nonnegative. Hence $\ell_B$ is nef.

A nef line bundle on an abelian variety whose associated symmetric
homomorphism is an isogeny is ample. To recall a short precise version
of this criterion, clear denominators and call the resulting bundle
$L$. Nondegeneracy gives $\chi(L)\ne0$ and
$\deg\phi_L=\chi(L)^2$, whereas nefness and Riemann--Roch give
$\chi(L)>0$. For an ample $H_0$, the polynomial
$\chi(L+tH_0)=(c_1(L)+tc_1(H_0))^{\dim B}/(\dim B)!$
is positive for all $t\ge0$. Mumford's index theorem therefore gives
index zero; its effective nondegenerate criterion then gives
ampleness \cite[Chapter~III, \S16, Index Theorem, and Chapter~II,
\S6, Application~1]{MumfordAV}. This argument is valid in arbitrary
characteristic. Applying it to \eqref{eq:T4-divisor} proves the theorem.
The zero-dimensional $B$ case is immediate. Rational coefficients
justify all divisions even when the characteristic divides a chosen
integral denominator; no assertion about an integral principal
polarization follows.
\end{proof}

\paragraph{Archive meaning and the exact input.}
This is the primitive reading required by the marginal instruction in
cote~44, scan~8, and the incidence product in scan~9. Cote~16,
scans~34, 36, 38, 40 supplies the Lefschetz language; scans~49--50 places
positivity in the ring of correspondences. The passage from cycle
positivity to an ample divisor is the curve test above; it is not the
unproved assertion that a cohomologically nondegenerate correspondence
is automatically a polarization. The input 032 is used only on
$J(C)\times A$ in Lemma~\ref{lem:curve-test}. No Hodge conjecture for a
non-CM characteristic-zero abelian variety and no new algebraicity
claim for such Hodge classes is inferred.

\paragraph{Sign calibration.}
For an elliptic curve, the normalized divisor
$P=\Delta-E\times\{0\}-\{0\}\times E$ induces
$x\mapsto\mathcal O_E(x-0)=-\phi_{\mathcal O_E(0)}(x)$.
Its self-incidence is $P$, and $\deg\Delta^*P=-2$.
Thus the sign in \eqref{eq:T4-divisor} is indeed $(-1)^m$ under the
stated convention: for $m=1$ it gives the positive elliptic
polarization. The frequently written sign $(-1)^{m-1}$ belongs instead
to the alternating primitive cohomological pairing before the
incidence-to-divisor identification. Confusing these two conventions
would reverse the conclusion already for a curve.
